Learned tactic policies are a core component of modern neural theorem provers: a language model proposes the next proof step and a best-first or tree search decides which proof states to expand~\cite{polu2020generative,lample2022hypertree,xin2025bfs,yang2023leandojo}. Earlier work used neural networks to rank premises or clauses for classical provers~\cite{alemi2016deepmath,loos2017deep} or to guide connection-tableau search with reinforcement learning~\cite{kaliszyk2018reinforcement}, and incremental learning with hindsight relabelling has been used to learn provers from few proofs~\cite{aygn2021proving}. In these systems, the learned component is evaluated on large libraries where a hand-written alternative is hard to build and optimal proofs are unknown, so it is difficult to say how much of a prover's success is due to the learned ranking, to the search algorithm and cost function, or to the problem distribution.

Propositional sequent calculus offers a controlled microcosm. With invertible (G3-style) rules, every tactic preserves provability, so a prover never gets stuck; what a good tactic choice buys is a \emph{shorter proof and less search}. Optimal proof size is computable for small sequents, a sound hand-written heuristic is known in principle (apply non-branching rules first so that context is not duplicated across branches), and sequents of arbitrary size are easy to generate. This makes it possible to ask a sharper question than ``does the policy help?'':

\emph{On sequents larger than anything in training, and under a fixed budget of \rhval{const/budget} expansions, does best-first search guided by a learned tactic policy find more proofs than breadth-first search and than a fairly tuned hand heuristic?}

We pre-registered four hypotheses (Sec.~\ref{sec:setup}): the policy beats BFS (H1) and the heuristic at the largest size (H2), the policy's relative standing worsens with size (H3), and search beats greedy decoding of the policy (H4). Our results are mixed to negative. H1 is supported; H2, H3 and H4 are refuted. An exploratory diagnostic and ablations indicate that how the policy's output is turned into a search priority matters more than the policy architecture, and that a feature-based MLP with a rank-based cost can exceed the heuristic in our setting, a post-hoc finding that we flag as such.

\paragraph{Contributions.} (i) A small, fully reproducible sequent-calculus prover, four random-provable-sequent generators and an optimal-proof labeller; (ii) a controlled comparison of BFS, a tuned hand heuristic, a feature MLP policy and a transformer policy under one search engine, with \rhval{const/n_seeds} seeds and three test size bins; (iii) ablations over the cost function, decoding mode, positional encoding, training-set size and budget, plus a diagnostic of how often the policies defer branching rules. We make no claim beyond this synthetic setting.
